%%% ========== 主文件导言区配置 ========== 

%% ======== 加载宏包 ========
\usepackage{subfiles}
\usepackage{lipsum,zhlipsum}
\usepackage{Macros/Ebox}
\usepackage{Macros/Ecodecopy}
\usepackage{Macros/Eenum}
\usepackage{Macros/Efiguretable}
%\usepackage{Macros/Emath}
%% ======== 宏包加载完毕 ======== 

%% ======== 封面配置 ======== 
% ------ 标题 ------
\title{VividBook MOD}
\subtitle{VividBook Modified Version}

% ------ 写作信息 ------
\author{3322we (大史莱姆)}
\date{\today}
\version{第 $2$ 版}
\bioinfo{邮箱}{tarquedumonde@gmail.com}

% ------ 寄语 ------
\extrainfo{一套功能超多, 报错又少的漂亮模板!}

% ------ 封面图片 ------
\logo{kling.png}
\cover{thecover.png}

% ------ 修改标题页的颜色带 ------
%\definecolor{customcolor}{RGB}{245, 250, 246}
\colorlet{coverlinecolor}{second!70}
%% ======== 封面配置完毕 ========

%% ======== 命令配置 ========
% ------ 公式控制 ------
\newcommand{\displaynarrow}{%
    \setlength{\abovedisplayskip}{3pt}
    \setlength{\belowdisplayskip}{3pt}%
    }
\newcommand{\displaywide}{%
    \setlength{\abovedisplayskip}{6pt}
    \setlength{\belowdisplayskip}{6pt}%
    }

% ------ 显示层级 ------
\setcounter{tocdepth}{2} %目录显示到 subsection 层级 
\setcounter{secnumdepth}{1} %标题编号只显示到 section 层级, 每节第一个子标题前面要加上命令 \phantomsection.

% ------ 自定义高亮 ------
\renewcommand{\emph}[1]{{\color{structurecolor}\bftitle#1}}
\newcommand{\fullemph}[1]{%
\colorbox{structurecolor!20}{#1}%
}

% ------ 允许环境分页显示 ------
\allowdisplaybreaks

% ------ 索引 ------
\usepackage{imakeidx}
\newcommand{\myindexheading}[1]{%
  \par\noindent\makebox[\linewidth][c]{%
    \color{structurecolor}\Large\bfseries #1%
  }\par\nopagebreak
}
\makeindex[options = -s indexstyle.ist]
%% ======== 命令配置完毕 ========
%%% ========== 导言区配置完毕 ========== 